Áno — jadro zadania vyzerá byť splnené, ale formálne to ešte treba dotiahnuť, aby to bolo obhájiteľné bez výhrad.

Podľa Anotacia.tex:1-8 máš 4 body:

1. Štúdia procedúr a algoritmov
Toto máš pravdepodobne splnené, lebo v main.tex:185-192 máš samostatné časti o nanopóroch, Dorade, clusteringu, Cloveri a multi-alignment.

2. Vytvorenie programových prostriedkov
Toto máš skôr splnené, keďže v main.tex:194-197 opisuješ vlastné spracovanie dát a hľadanie v DNA reťazcoch.
Tu však odporúčam explicitne pomenovať, čo presne si vytvoril:

aké skripty/nástroje,
aký majú vstup,
aký majú výstup,
na čo sa používajú.
3. Dokumentácia programových prostriedkov
Toto je zatiaľ najslabší bod. V štruktúre práce to nie je úplne viditeľné ako samostatná dokumentácia.
Odporúčané doplniť:

stručný popis každého nástroja,
spôsob spustenia,
vstupy/výstupy,
použité parametre,
príklad použitia,
obmedzenia riešenia.
4. Kritické zhodnotenie výsledkov
Toto máš zatiaľ iba veľmi stručne. Záver v main.tex:200-201 je príliš všeobecný.
Treba doplniť:

čo sa podarilo konkrétne,
čo sa nepodarilo alebo ostalo otvorené,
presnosť/limity/metriky, ak ich máš,
čo by bolo ďalším krokom.
Dôležité upozornenia
Abstrakt ešte nemáš hotový.
V abstract.tex:19-38 je stále placeholder text. Toto určite treba nahradiť reálnym abstraktom v slovenčine aj angličtine.

Pozor na konzistentnosť cieľa práce.
V NaplnUlohy.tex:3-14 opisuješ cieľ ako simuláciu chýb a generátor, čo nie je úplne totožné s formuláciou anotácie. Toto by som zjednotil, aby cieľ, anotácia, abstrakt a záver hovorili to isté.

Samotná anotácia je jazykovo slabšia.
V Anotacia.tex:1-8 sú chyby typu „algorimtmov“, chýbajú dĺžne a mäkčene. Ak je to text, ktorý bude niekde oficiálne použitý, treba ho preštylizovať.

Verdikt
Obsahovo: áno, smeruješ k splneniu zadania.
Formálne: ešte by som doplnil dokumentáciu nástrojov, silnejšie kritické zhodnotenie a hotový abstrakt.

